\chapter{Einleitung}
\label{cha:sa_Einleitung}

\section{Zielsetzung}
Dieses Dokument ist als vorwiegend technische Starthilfe für das
Erstellen einer Diplomarbeit mit \LaTeX
gedacht.

Das vorliegende Dokument dient einem zweifachen Zweck: 
\emph{erstens} als Erläuterung und Anleitung, \emph{zweitens} als
direkter Ausgangspunkt für die eigene Arbeit.


\section{Warum {\LaTeX}?}

Diplomarbeiten, Dissertationen und Bücher im
technisch-natur\-wissen\-schaft\-lichen Bereich werden
traditionell mithilfe des Textverarbeitungssystems \LaTeX
\cite{Lamport94,Lamport95} gesetzt.

Für den professionellen elektronischen Buchsatz wurde früher
häufig \emph{Adobe Framemaker} verwendet, allerdings ist diese
Software teuer und komplex. Eine modernere Alternative dazu ist
\emph{Adobe InDesign}, wobei allerdings die Erstellung
mathematischer Elemente und die Verwaltung von Literaturverweisen
zur Zeit nur rudimentär unterstützt werden.%
\footnote{Angeblich werden aber für den (sehr sauberen) Schriftsatz 
	in \emph{InDesign} ähnliche Algorithmen wie in \LaTeX\ verwendet.}

Microsoft \emph{Word} gilt im Unterschied zu \LaTeX, 
\emph{Framemaker} und \emph{InDesign} übrigens nicht als professionelle
Textverarbeitungssoftware, obwohl es immer häufiger auch von
großen Verlagen verwendet wird.%
\footnote{Siehe auch \url{http://latex.tugraz.at/mythen.php}.}
Das Schriftbild in \emph{Word}
lässt -- zumindest für das geschulte Auge -- einiges zu wünschen
übrig und das Erstellen von Büchern und ähnlich großen Dokumenten
wird nur unzureichend unterstützt. Allerdings ist \emph{Word} sehr
verbreitet, flexibel und vielen Benutzern zumindest oberflächlich
vertraut, sodass das Erlernen eines speziellen Werkzeugs wie
\LaTeX\ mühsam sein kann.

Übrigens, genau hier am Ende des Einleitungskapitels (und nicht
etwa in der Kurzfassung) ist der richtige Platz, um die
inhaltliche Gliederung der nachfolgenden Arbeit zu beschreiben.
Hier soll dargestellt werden, welche Teile (Kapitel) der Arbeit
welche Funktion haben und wie sie inhaltlich zusammenhängen. Auch
die Inhalte des \emph{Anhangs} -- sofern vorgesehen -- sollten hier
kurz beschrieben werden.

